% ====================================================================
% §18 전체 입력 인자와 피팅 준비 — 진행 요약
% ====================================================================
\section{전체 입력 인자와 피팅 준비 --- 진행 요약}\label{sec:inputs}
사용자 목표는 이 모든 인자를 입력으로 노출해 측정 데이터에 피팅하는 것이다. 곡선식의 모든 물리량이 사용자
조정 가능한 입력이며, 내부에 고정된 상수는 없다(전부 기본값$+$override --- 물리 상수 $R,F,k_B,h$ 와
수치 안전 가드 상수는 물리 인자가 아니므로 제외). 전 조정 인자와 기본값의 상세
목록은 부록~\ref{sec:appendix-code}(표~\ref{tab:inputs})에 정리한다 --- 어느 것을 고정하고 어느 것을
피팅할지는 사용자가 데이터로 정한다(본 장은 스코프를 정하지 않는다).

``고정$\cdot$피팅 스코프''는 자유롭게 잡을 수 있다 --- 예컨대
$\Delta H_\rxn$$\cdot$$\Delta S_\rxn$$\cdot$$Q$$\cdot$$w$ 를 풀고 $\Omega$$\cdot$$\Delta H_a$$\cdot$$|\dd V/\dd q|_{q_a}$
는 좁은 상하한으로 두는 식이다. 반응 열역학($\Delta H_\rxn\cdot\Delta S_\rxn$) 값과 DFT 활성화 값은 곧이
믿고 잠글 경계(상$\cdot$하한)가 아니라 \emph{출발점}이므로 소폭 자유도를 권한다.

\begin{warnbox}
\textbf{식별 가드 --- $\alpha$$\cdot$$L_V$$\cdot$$w$$\cdot$gallery 네 손잡이의 축퇴(v1.0.25).} 정적 단일
온도$\cdot$단일 율속 곡선 하나에서 봉우리의 \emph{비대칭}과 \emph{정점 이동}은 서로 다른 네 손잡이가
거의 같은 방식으로 만들 수 있다 --- 넷을 동시에 자유화하지 말 것.
\begin{enumerate}[label=(\arabic*),leftmargin=2.0em,itemsep=2pt]
\item $\alpha_j$ --- 평형 skew 지수(식~\eqref{eq:skewpeak}). 정점을 $\sigma_d w_j\ln\alpha_j$ 만큼 밀고
반높이 폭까지 함께 바꾼다($\alpha_j$ 와 함께 단조 감소 --- 식~\eqref{eq:skewapex} 문단).
\item $L_{V,j}$ --- 유한율속 인과 꼬리(\S\ref{sec:lag}). 정점을 진행방향으로 밀고 한쪽 꼬리를 늘인다.
\item $w_j$(곧 $n_j$) --- 폭 척도(식~\eqref{eq:wbase}). $\alpha_j$ 의 폭 효과와 직접 겹친다.
\item gallery 세분 --- 한 물리 봉우리를 중심이 거의 같고 폭이 다른 두 전이로 나누면 그 envelope 이 곧
비대칭 봉우리가 된다(중심 간격이 폭보다 작을 때).
\end{enumerate}
\emph{권장 스코프} --- $\alpha_j$ 를 열 때는 같은 전이의 $L_{V,j}$ 를 동결하고(구현은 같은 전이에
$\alpha\ne1$ 과 $L_V>0$ 이 동시에 지정되면 경고를 낸다), $\alpha_j$ 와 $n_j$ 중 하나만 푼다.
gallery 세분과 $\alpha_j$ 의 동시 자유화는 \emph{조건부}다(v1.0.25.2 조준 정정): 둘은 위 (1)$\cdot$(4)
로 축퇴하므로 \emph{개별 값}을 물리량으로 읽으려면 동시에 열지 말아야 하지만, 목적이 \emph{곡선 표현}
이라면 열어도 된다 --- 모수 보정 지표가 전이당 벌점을 물고도 두 손잡이를 함께 연 쪽을 큰 차로 택하므로,
이 데이터에서 둘은 완전 축퇴가 아니다. 곧 \textbf{동시 자유화의 대가는 적합도가 아니라 개별 $\alpha_j$
$\cdot$$w_j$ 의 점-식별성}이며, 그렇게 얻은 값은 seed 이지 신뢰값이 아니다(불확실도 구간을 함께 내지
않으면 인용 금지). 어느 경우에도 $L_{V,j}$ 동결은 유지한다 --- (2)만이 율속 스윕으로 분리되는 유일한
몫이라, 그것까지 함께 열면 남은 셋의 분리 근거가 사라진다.
\emph{원리적 분리} --- 넷 중 (2)만 율속에 비례해 자라고 $|I|\!\to\!0$ 에서 소멸하므로
(\S\ref{sec:broadening-sources} ①), \emph{율속 스윕}이 유일한 분리 관측이다. (1)$\cdot$(3)$\cdot$(4)의
분리는 단일 곡선으로는 비유일하며, 다온도$\cdot$다율속 데이터 또는 독립 구조 정보(XRD 상 수 ---
\S\ref{sec:broadening-class})가 있어야 한다. 그러므로 이 손잡이들로 얻은 값은 \emph{곡선 표현}의
파라미터이지 그 자체로 상(phase) 개수$\cdot$상 성격의 측정이 아니다.
\emph{다섯째 손잡이 --- $\Omega$(v1.0.25.2).} 두-상 파라미터 $\Omega$ 를 곡선 적합에서 함께 풀면, 그것은
$\Omega\!\to\!2RT^+$ 에서 봉우리를 임의로 좁힐 수 있으므로 (1)$\cdot$(3) 과 같은 자리에서 축퇴한다. 실측
징후는 명확하다 --- 적합 결과의 $\Omega$ 가 탐색 구간의 \emph{경계에 붙으면} 그 값은 식별된 것이 아니다
(경계 위 어느 값이어도 같은 곡선이 나온다). 그런 $\Omega$ 는 두-상 판정(\S\ref{ssec:gr-2L})의 근거로 쓰지
말 것 --- 물리량이 아니라 폭 손잡이로 쓰인 것이다.
\end{warnbox}

\begin{warnbox}
\textbf{적합도 이전에 볼 것 --- 전하 보존(v1.0.25.2).}

\emph{(가) 면적 보존은 적합도와 \emph{독립된} 진단이다.} 본 장의 종은 전부 면적이 $Q_j$ 로 닫히도록
정규화되어 있으므로(식~\eqref{eq:eqpeak}$\cdot$\eqref{eq:skewpeak}), 적합 창 안에서
\[
\Delta_Q\;\equiv\;1-\frac{\displaystyle\int_\text{창}\Big(\frac{\dd Q}{\dd V}\Big)^\text{모델}\dd V}
{\displaystyle\int_\text{창}\Big(\frac{\dd Q}{\dd V}\Big)^\text{실측}\dd V}
\]
는 0 에 가까워야 한다. $\Delta_Q$ 는 잔차 제곱합$\cdot$정보 기준과 \emph{다른 축}이다 --- 결정계수가
높아도 급준한 봉우리를 세우지 못해 창 안 전하를 수 \%\ 놓치는 적합이 실제로 나온다. 그러므로 적합을
채택하기 전에 $\Delta_Q$ 를 먼저 보고, 크면 (결정계수와 무관하게) 그 적합은 \emph{전하 보존을 깨뜨린
것}으로 판정한다. 본 문건의 전체 골격이 전하 보존식 위에 서 있으므로 이 진단은 선택이 아니다.
\emph{주의}: 배경 항이 자유롭게 열려 있으면 과대한 배경이 놓친 봉우리 면적을 메워 $\Delta_Q$ 를 좋게
보이게 할 수 있다 --- $\Delta_Q$ 는 배경 값의 타당성 점검과 \emph{함께} 읽어야 한다.

\emph{(나) 실측 곡선 쪽 전제.} 위 진단이 성립하려면 실측 $\dd Q/\dd V$ 자체가 전위축에서 균일하게
표본된 값이어야 한다 --- 계측 전위는 유한 눈금으로 양자화되어 있어 평탄역에서 표본 간 전위차가 눈금 이하로
내려가므로, 미분을 그대로 취하면 하필 \emph{두-상 평탄}에서 발산한다. 곧 이 발산은 물리가 아니라 표본화의
산물이다. 실측 곡선을 만드는 수치 절차는 물리식이 아니므로 부록~\ref{sec:appendix-code} 에 둔다.
\end{warnbox}

\begin{keybox}
\textbf{본 장만으로 한 곡선 재현(6단계).}
(1) 전이 파라미터 집합을 구성 --- 기준선은 7-gallery skew($\alpha_j$ 포함, \S\ref{ssec:gr-2L})이고,
표~\ref{tab:staging} 의 4 전이는 선택(상 수와 1:1 대응하는 물리 기준$\cdot$회귀 비교용)이다.
(2) 실험조건 $\sigma_d$, $|I|=$ c-rate$\,\cdot Q_\cell$, $T$, $R_n$ 설정 → 분극 $V_n$(\eqref{eq:n0map}\eqref{eq:vn}); 이후 평가는 $V_n$ 에서 점별.
(3) 전이마다 $U_j(T)$(\eqref{eq:Uj}) → 분기 중심 $U_j^{\,d}$(\eqref{eq:center}) → 폭 $w_j$(\eqref{eq:wbase}) → 평형 진행률 $\xi_\eq$(\eqref{eq:xieq}).
(4) 지연 길이: $\mathcal A$(\eqref{eq:Acut}) → $\chi_d,\Delta H_a^\eff$(\eqref{eq:chid}\eqref{eq:dHeff}) → $L_q$(\eqref{eq:Lqfull}) → $L_V$(\eqref{eq:LV}), 전이당 상수.
(5) peak 모양: 기억 적분(\eqref{eq:lag}$\cdot$\eqref{eq:reversal})의 $(\xi_\eq-\xi_\mathrm{lag})/L_V$(\eqref{eq:peakshape}); $L_V\to0$ 극한이 평형 종(\eqref{eq:tail-limit}\eqref{eq:eqpeak}).
(6) 용량 가중 합산(\eqref{eq:sum}); 출력이 입력 전위 $V_n$ 좌표에서 곧바로 나온다. 색인은 표~\ref{tab:nodemap}.
\end{keybox}

\subsection{전체 진행 요약}\label{sec:facade}
요약하면 --- 계산은 실험조건을 받아 방향 부호와 전류 크기로 먼저 환산하고(식~\eqref{eq:n0map}),
LCO 전극은 셀 라벨을 탈리튬화 부호로 자동 환산하며(식~\eqref{eq:lco-sigmaslot}), $|I|\to0$ 기준선은
식~\eqref{eq:eqpeak} 만 합산한 충방전 불변$\cdot$히스테리시스 미반영 곡선이다. 전체 진행은
표~\ref{tab:nodemap} 가 노드와 식으로 정리하고, 구현 대응표는
부록~\ref{sec:appendix-code} 에 둔다.

\begin{table}[t]
\centering
\renewcommand{\arraystretch}{1.3}
\caption{노드 $\leftrightarrow$ 식 매핑(전체 진행). 본 장 절 순서 $=$ 이 노드 순서.
구현 식별자 대응은 부록~\ref{sec:appendix-code}(표~\ref{tab:nodecode}).}
\label{tab:nodemap}
\begin{tabular}{l l l}
\toprule
노드 & 물리식 & 본 장 식 \\
\midrule
N0 & $\sigma_d$, $|I|=$ c-rate$\,\cdot Q_\cell$ & \eqref{eq:n0map} \\
N1 & $V_n=V_\app-\sigma_d|I|R_n$ & \eqref{eq:vn} \\
N2 & $U_j(T)=(-\Delta H+T\Delta S)/F$ & \eqref{eq:Uj} \\
N3 & $\Delta U_j^\hys$, $U_j^{\,d}$ & \eqref{eq:dUhys},\eqref{eq:Ubranch} \\
N4 & $w_j=n_jRT/F$ (이중지위) & \eqref{eq:wbase} \\
N5 & $\xi_\eq=\mathrm{logistic}[\sigma_d(V-U^d)/w]$ & \eqref{eq:xieq} \\
N6 & $Q_j\xi_\eq(1-\xi_\eq)/w$ & \eqref{eq:eqpeak} \\
\multicolumn{3}{l}{\emph{\quad(N6 확장: 두-상 broadening$\cdot$현상학적 $w_j$$\cdot$역산 금지 = \S\ref{sec:broadening}; 새 N-노드 아님, 모델 항 0)}}\\
N7 & $\mathcal A,\chi_d,\Delta H_a^\eff,L_q,L_V$ & \eqref{eq:Acut}--\eqref{eq:LV} \\
N8 & $(\xi_\eq-\xi_\mathrm{lag})/L_V$, 역전 & \eqref{eq:peakshape},\eqref{eq:reversal} \\
N9 & $C_\bg+\sum_jQ_j[\cdot]$ & \eqref{eq:sum} \\
\midrule
\multicolumn{3}{l}{\emph{— LCO 양극 추가(같은 골격, 파라미터+1항) —}}\\
N0$'$ & LCO 전이 초기값(3 전이) & 표~\ref{tab:lco-staging} \\
\multicolumn{3}{l}{\emph{\quad(N0$'$ 보강: 셀 라벨$\to$탈리튬화 부호 자동 환산 = 식~\eqref{eq:lco-sigmaslot}$\cdot$\S\ref{sec:facade})}}\\
N2$'$ & $\partial U_j/\partial T=\Delta S^\mathrm{cat}/F$ & \eqref{eq:lco-dUdT} \\
N3$'$ & $\Delta U_j^{\hys,\mathrm{cat}}$, $U_j^{\,d,\mathrm{cat}}$ & \eqref{eq:lco-dUhys},\eqref{eq:lco-Ubranch} \\
N5$+$ & $\Delta S_{e,j}^{\,\mathrm{mol}}=N_A\tfrac{\pi^2}{3}k_B^2T\,\partial g/\partial x<0$ (삽입, 몰당) & \eqref{eq:dSemolar},\eqref{eq:ggate} \\
N9$'$ & $\Delta S^\mathrm{cat}=$config$+$vib$+$elec & \eqref{eq:lco-decomp} \\
\bottomrule
\end{tabular}
\end{table}

여기까지가 본 장(흑연)의 결과 사슬이다 --- 표의 LCO 행(N$'$ 계열)은 Chapter 2 가 같은 forward 골격을
재사용하는 지점의 통합 대응 참조다. 부록 A 가 부호 사슬을
전수 검산하고, 부록 B 가 구현과의 1:1 대응을 역방향으로 조회한다.

% 자산: [B2-127] [B2-128] [B2-129]
